\documentclass[a4paper,12pt]{article}
\usepackage{amsmath, amssymb}
\usepackage{geometry}
\usepackage{enumitem}
\usepackage{fancyhdr}
\usepackage{xcolor}
\usepackage{sectsty}
\usepackage{multicol}
\usepackage{graphicx}
\usepackage{float}
\def\inputGnumericTable{}
\usepackage[latin1]{inputenc}
\usepackage{fullpage}
\usepackage{color}
\usepackage{array}
\usepackage{longtable}
\usepackage{calc}
\usepackage{multirow}
\usepackage{hhline}
\usepackage{ifthen}

\geometry{top=1in, bottom=1.5cm, left=1.5cm, right=1.5cm}
\pagestyle{empty}
\sectionfont{\color{blue}}

\begin{document}

\thispagestyle{fancy}
\fancyhf{}
\fancyhead[L]{
        
\includegraphics[width=8cm, height=1.7cm]{images/logo.png}
}
\fancyhead[R]{
    Name: Bhoomika B V \\
    Batch: COMETFWC046 \\
    Date: \today
}
\renewcommand{\headrulewidth}{0pt}
\fancyfoot[C]{\thepage}

\vspace*{1cm}

\begin{center}
	{\LARGE \textbf{\textcolor{blue}{GATE EC Question Analysis using Seven Segment Display}}}
\end{center}

\section*{\textbf{Question}}
If $X = 1$ in the logic expression

\[
[X + Z(\bar Y + (\bar Z + X\bar Y))][\bar X + \bar Z(X+Y)] = 1
\]

Find the condition for output to be 1.

\section*{\textbf{Question Analysis:}}

Substitute $X = 1$

\[
[1 + Z(\bar Y + (\bar Z + \bar Y))][0 + \bar Z]
\]

\[
= 1 \cdot \bar Z
\]

\[
\therefore F = \bar Z
\]


\section*{\textbf{The Truth Table}}

\begin{center}
\begin{tabular}{|c|c|}
\hline
Z & F \\ \hline
0 & 1 \\
1 & 0 \\
\hline
\end{tabular}
\end{center}

\section*{\textbf{Hardware Implementation}}

The above problem is implemented and tested in hardware using Arduino UNO board.  
A seven segment display is used to display output $F$ as decimal digit 0 or 1 based on input $Z$.  
The Arduino reads input $Z$, computes $F=\bar Z$ and displays the output.

\section*{Required Components \& Pin Connections}

\begin{center}
\begin{minipage}{0.45\textwidth}
\begin{table}[H]
\centering
\begin{tabular}{|c|l|}
\hline
\textbf{S.No} & \textbf{Component} \\ \hline
1 & Arduino Uno Board \\
2 & Breadboard \\
3 & Seven segment display \\
4 & Resistors: 220$\Omega$ \\
5 & Jumper Wires \\
6 & USB Cable \\
\hline
\end{tabular}
\end{table}
\end{minipage}
\hspace{0.05\textwidth}
\begin{minipage}{0.45\textwidth}
\begin{table}[H]
\centering
\begin{tabular}{|c|c|}
\hline
\textbf{Component} & \textbf{Arduino Pin} \\ \hline
Input Z & Digital 9 \\
Segment a & Digital 2 \\
Segment b & Digital 3 \\
Segment c & Digital 4 \\
Segment d & Digital 5 \\
Segment e & Digital 6 \\
Segment f & Digital 7 \\
Segment g & Digital 8 \\
GND & GND \\
VCC & 5V \\
\hline
\end{tabular}
\end{table}
\end{minipage}
\end{center}

\section*{Logic Description}

\begin{itemize}
	\item Read input $Z$
	\item Simplified Boolean expression: $F=\bar Z$
	\item If $Z=0$ display digit 1
	\item If $Z=1$ display digit 0
\end{itemize}

\section*{Code Uploading Steps}

\begin{enumerate}
	\item Create Platform IO project
	\item Write code in main.cpp
	\item Compile using command ``pio run''
	\item Connect Arduino UNO using USB
	\item Upload compiled code
	\item Observe output on seven segment display
\end{enumerate}

\section*{Experimental Truth Table}

\begin{center}
\begin{tabular}{|c|c|}
\hline
Z & Display Output \\ \hline
0 & 1 \\
1 & 0 \\
\hline
\end{tabular}
\end{center}

\begin{center}
	\includegraphics[width=12cm, height=12cm]{images/output1.jpg}
\end{center}

\section*{Conclusion}

\begin{itemize}
    	\item The Boolean expression simplifies to $F=\bar Z$
	\item Hardware output matches theoretical truth table
    	\item Seven segment display successfully verifies the logic
\end{itemize}

\end{document}

